\scanpage{25}
\mspage{24}
\begin{center}-4.14-\end{center}

of $\overline V_4$ are precisely the four lines $\ell_{ij}\subset\overline V_i\subset\overline V_4$, and $H_4$ passes through all of them. A local computation shows that these singularities are all resolved on $\overline V_5$, and each line $\ell_{ij}$ is blown up into a ruled surface $R_{ij}$. Moreover, the action of $\mathbb G_m$ lifts to $\overline V_5$, since $H_4$ is invariant under $\mathbb G_m$, and $\mathbb G_m$ acts non-trivially on the ruled surface $R_{ij}$ with orbit space $\ell_{ij}$. Finally, define $V_5$ to be the opensubset of $\overline V_5$ obtained by removing:
\margin{$\unsure{\sim\sim}$}

a) the proper transforms of $\overline V_1$, $\overline V_2$ ;

b) the proper transforms on $\overline V_5$ of the divisors $E_1$ and $E_2$ on $\overline V_4$; where $E_1$ is the locus of lines joining $\ell_{11}$ and $\ell_{21}$, and $E_2$ is the locus of lines joining $\ell_{12}$ and $\ell_{22}$ ;

c) $R_{12}$ and $R_{21}$.

A rather conceptualized ``picture'' of $V_5$ is given in Figure 8.

It is clear that $V_5$ is invariant under the action of $\mathbb G_m$, and that $V_3$ is a geometric quotient of $V_5$ by $\mathbb G_m$: for $V_5$ is the union of the lines joining biregularly corresponding points of $V_1$ and $V_2$, plus the lines of $R_{11}$ and of $R_{22}$.

Note also that the automorphism $\sigma_{\unsure{o|0}}$ can be carried along:

i) $\sigma$ lifts to an isomorphism $\overline V_1\longrightarrow\overline V_2$ ; ii) if the embeddings $\overline V_i\subset\mathbb P_{n_i}$ are chosen suitably, $\sigma$ extends to an isomorphism $\mathbb P_{n_1}\longrightarrow\mathbb P_{n_2}$ ; iii) this isomorphism extends to an automorphism of $\mathbb P_{n_1+n_2+1}$ of order 2 satisfying $\sigma.\alpha.\sigma^{-1}=\unsure{-}\alpha^{-1}$, for $\alpha\in(\unsure{\overline{\mathbb G}_m|\mathbb G_m})$ acting on $\mathbb P_{n_1+n_2+1}$ ; iv) $\sigma$ takes $\overline V_4$ into itself, and $H_4$ into itself, hence $\sigma$ lifts to a morphism $\sigma:\overline V_5\longrightarrow\overline V_5$; v) this $\sigma$ carries
